% ============================================================================
% Fichier  : partie2/P02_C04_analyse_disque.tex
% Titre    : L'Analyse du Disque
% Auteur   : MINKA MI NGUIDJOÏ Thierry Emmanuel
% Version  : 1.0 -- novembre 2026
% Corrigés : corriges/Corr_C04.tex
% ============================================================================
% Sources : Carrier (2005) pour les systèmes de fichiers ; documentation
% Microsoft des journaux de sécurité (identifiants 4624, 4625, 4688, 4720,
% 7045, 1102, présentés sans citation) ; The Sleuth Kit et Autopsy.
% Faits nouveaux (bible) : sur SRV-CPT01, session RDP de adm.kamga depuis
% SRV-AD01 le 12/03 à 22 h 39 ; raccourci api.ini.lnk créé à 22 h 44 ;
% archive cpt.7z créée puis supprimée à 22 h 52, récupérée ; prefetch
% désactivé (Windows Server) ; journal Amcache.
% ============================================================================

\chapter{L'Analyse du Disque}
\label{chap:disque}
\chaptermark{L'analyse du disque}

\epigraph{%
    « Supprimer un fichier, pour un système de fichiers,\\
    c'est d'abord oublier où il se trouvait. »%
}{Principe du cours}

\niveauChapitre{Intermédiaire}{%
    Maîtrise du Chapitre~\ref{chap:acquisition}
}

\begin{boxobjectifs}
\begin{itemize}
    \item Lire la structure d'un disque~: table de partitions, volumes,
          système de fichiers.
    \item Expliquer comment NTFS enregistre un fichier, ses dates et sa
          suppression.
    \item Récupérer un fichier supprimé par ses métadonnées ou par ses
          signatures.
    \item Exploiter les artefacts Windows qui prouvent une connexion, une
          exécution ou l'ouverture d'un fichier.
    \item Interpréter un horodatage avec son fuseau et ses limites.
\end{itemize}
\end{boxobjectifs}

\bigskip

% ============================================================================
\section*{Avant de commencer : la question posée au disque}
% ============================================================================

Yaoundé, laboratoire de SANKOFA EXPERTISE, 19 mars 2025. La copie de travail
de l'image de SRV-CPT01 est montée en lecture seule. Le disque compte un peu
plus de 300\,000 fichiers. Mme~MBARGA n'ouvre pas l'outil d'analyse au
hasard~: elle écrit d'abord trois questions. L'attaquant s'est-il connecté à
ce serveur, et quand~? A-t-il ouvert le fichier qui contient la clé de
l'interface de paiement~? A-t-il copié quelque chose et effacé ses traces~?

\begin{boxalert}
\textbf{Violation typique.} Parcourir l'arborescence d'une image en
espérant « tomber sur quelque chose », puis présenter ce que l'on a trouvé
comme le résultat d'un examen complet.

\smallskip
\textbf{Règle~:} on examine un disque pour répondre à des questions écrites
d'avance, en choisissant les artefacts capables d'y répondre, et l'on dit
pour chaque question quelles sources ont été lues.
\end{boxalert}

\separateur

% ============================================================================
\section{Du Disque au Fichier}
\label{sec:c04-structure}
% ============================================================================

Un disque se lit par couches. La \termecle{table de partitions} (MBR sur les
anciens systèmes, GPT sur les récents) découpe le disque en volumes. Chaque
volume porte un \termecle{système de fichiers}, qui organise les données en
fichiers et en répertoires~: NTFS sous Windows, ext4 sous la plupart des
Linux, APFS sous macOS. L'espace d'un volume se répartit entre les données
allouées, les métadonnées du système de fichiers, et l'espace non alloué,
qui contient souvent les restes de fichiers supprimés.

\begin{center}
\begin{tikzpicture}[
    couche/.style = {draw=CouleurBordDef, fill=CouleurDef, minimum width=10cm,
        minimum height=0.75cm, align=center, font=\small\color{CouleurPrimaire}}
]
\node[couche] at (0,0)    {Disque physique~: secteurs de 512 ou 4\,096 octets};
\node[couche] at (0,0.9)  {Table de partitions (MBR, GPT)~: où commence chaque volume};
\node[couche] at (0,1.8)  {Système de fichiers (NTFS, ext4)~: métadonnées et données};
\node[couche] at (0,2.7)  {Fichiers et répertoires~: ce que voit l'utilisateur};
\node[couche] at (0,3.6)  {Artefacts~: journaux, registre, raccourcis, traces d'exécution};
\end{tikzpicture}
\end{center}
\begin{center}{\small\itshape Figure~4.1. Les couches de lecture d'un disque}\end{center}

% ============================================================================
\section{NTFS : la Table MFT et les Dates}
\label{sec:c04-ntfs}
% ============================================================================

NTFS décrit chaque fichier dans un enregistrement de la \termecle{table
MFT} (\textit{Master File Table}). Un enregistrement contient des
attributs~: le nom, les dates, la taille et l'emplacement des données sur le
disque. Les petits fichiers, de quelques centaines d'octets, tiennent même
entièrement dans leur enregistrement.

\subsection{Deux jeux de quatre dates}

Chaque fichier porte quatre dates~: modification du contenu (M), dernier
accès (A), modification de l'enregistrement MFT (C), et création (B). On les
note MACB. NTFS les enregistre deux fois~: dans l'attribut
\texttt{\$STANDARD\_INFORMATION}, que les programmes modifient librement,
et dans l'attribut \texttt{\$FILE\_NAME}, que seul le système met à jour lors
de certaines opérations (création, renommage, déplacement). Un écart
incohérent entre les deux jeux est l'un des signes d'une falsification des
dates, étudiée au Chapitre~\ref{chap:antiforensique}.

\begin{boxinfo}
\textbf{La date d'accès ne vaut pas grand-chose.} Les versions récentes de
Windows ne mettent pas toujours à jour la date de dernier accès, pour des
raisons de performance, et un antivirus qui lit un fichier la modifie. On ne
fonde jamais une conclusion « l'attaquant a ouvert ce fichier » sur la seule
date d'accès.
\end{boxinfo}

\subsection{Les journaux de NTFS}

Le journal \texttt{\$UsnJrnl} enregistre les changements des fichiers
(création, renommage, suppression), avec leur heure, tant que sa taille le
permet. Il permet de savoir qu'un fichier a existé, sous quel nom, et quand
il a été supprimé, même si son enregistrement MFT a été réutilisé.

% ============================================================================
\section{Les Fichiers Supprimés}
\label{sec:c04-supprimes}
% ============================================================================

Supprimer un fichier, sous NTFS, revient à marquer son enregistrement MFT et
ses blocs comme libres. Les données restent sur le disque jusqu'à ce qu'un
autre fichier les écrase. Deux méthodes permettent de les retrouver.

\begin{tableaucours}{L{3.2cm}L{4.6cm}L{4.4cm}}{%
    \tableauHeader{Méthode} &
    \tableauHeader{Principe} &
    \tableauHeader{Limite}
}
Par les métadonnées & L'enregistrement MFT libre indique encore le nom, les
dates et l'emplacement des données & Échoue si l'enregistrement a été
réutilisé \\
Par les signatures (\textit{carving}) & On cherche dans l'espace non alloué
les octets qui marquent le début d'un format (\texttt{PK} pour une archive
ZIP, \texttt{7z} pour une archive 7-Zip, \texttt{\%PDF} pour un PDF) & Ni
nom ni date~; fichiers fragmentés mal reconstitués \\
\end{tableaucours}
\vspace{-0.8em}
\begin{center}{\small\itshape Tableau~4.1. Deux façons de récupérer un fichier}\end{center}

Sur un disque SSD, la commande TRIM réduit fortement les chances de
récupération (Chapitre~\ref{chap:acquisition}). Sur un serveur virtualisé, le
comportement dépend de l'hyperviseur.

% ============================================================================
\section{Les Artefacts Windows}
\label{sec:c04-artefacts}
% ============================================================================

Un \termecle{artefact} est une trace qu'un système laisse en fonctionnant
normalement, sans qu'on lui demande de garder une preuve~: une entrée de
journal, une clé de registre, un raccourci. L'enquêteur les choisit d'après
la question.

\begin{tableaucours}{L{3.0cm}L{4.4cm}L{4.8cm}}{%
    \tableauHeader{Question} &
    \tableauHeader{Artefact} &
    \tableauHeader{Ce qu'il dit}
}
Qui s'est connecté, quand, comment~? & Journal de sécurité~: événements
4624 (ouverture réussie), 4625 (échec), type d'ouverture 10 pour le bureau
à distance & Compte, heure, adresse d'origine \\
Quel programme a été exécuté~? & Événement 4688 si l'audit est activé~;
Amcache~; Prefetch sur les postes de travail & Nom, chemin, heure, parfois
empreinte \\
Quel service a été installé~? & Événement 7045 du journal système & Nom du
service, binaire, compte \\
Quel fichier a été ouvert~? & Raccourcis \texttt{.lnk} du dossier
\texttt{Recent}, listes de raccourcis, clés \texttt{RecentDocs} du
registre de l'utilisateur & Chemin du fichier, heure d'ouverture \\
Le journal a-t-il été effacé~? & Événement 1102 & Effacement du journal
de sécurité, compte qui l'a fait \\
Un compte a-t-il été créé~? & Événement 4720 & Nom du compte, créateur \\
\end{tableaucours}
\vspace{-0.8em}
\begin{center}{\small\itshape Tableau~4.2. Questions et artefacts Windows}\end{center}

Deux précautions accompagnent ce tableau. Les artefacts dépendent de la
configuration~: le Prefetch, utile sur un poste de travail, est désactivé
par défaut sur les versions serveur de Windows, et l'événement 4688 n'existe
que si l'audit de la création de processus a été activé. Et chaque artefact
dit moins qu'il ne paraît~: une entrée Amcache prouve qu'un programme était
présent et a été vu par le système, pas toujours qu'il a été exécuté à
l'heure indiquée. On croise donc les artefacts entre eux.

% ============================================================================
\section{Ce que dit le Disque de SRV-CPT01}
\label{sec:c04-nyanga}
% ============================================================================

\begin{boxscenario}{Opération NYANGA. Trois questions, trois réponses}
\textbf{Question 1~: connexion.} Le journal de sécurité contient, le
12~mars 2025 à 22~h~39, un événement 4624 de type 10 pour le compte
\texttt{adm.kamga}, depuis l'adresse interne de SRV-AD01. Aucune ouverture
de ce compte sur SRV-CPT01 n'existe dans les trente jours précédents.

\textbf{Question 2~: ouverture du fichier de clé.} Le dossier \texttt{Recent}
du profil \texttt{adm.kamga} contient un raccourci \texttt{api.ini.lnk},
créé le 12~mars à 22~h~44, qui pointe vers
\texttt{C:\textbackslash PaySuite\textbackslash conf\textbackslash api.ini}.
La clé \texttt{RecentDocs} du registre de l'utilisateur cite le même
fichier.

\textbf{Question 3~: copie et effacement.} Le journal \texttt{\$UsnJrnl}
enregistre la création, à 22~h~52, d'un fichier \texttt{cpt.7z} dans le
dossier temporaire de l'utilisateur, puis sa suppression à 23~h~06.
L'enregistrement MFT est libre mais intact~: l'archive est récupérée, 41~Mo,
protégée par un mot de passe. Sa liste de fichiers, lisible sans le mot de
passe, comprend \texttt{api.ini} et le dossier des exports comptables de
février.

\medskip
\textbf{Analyse} (confiance élevée). Le compte \texttt{adm.kamga} a été
utilisé depuis le contrôleur de domaine pour ouvrir le fichier de clé, puis
pour l'archiver avec des exports comptables, avant d'effacer l'archive. Le
journal de sécurité ne montre aucun effacement (pas d'événement 1102)~:
l'attaquant a supprimé son archive, pas les journaux.
\end{boxscenario}

\begin{boxobjet}{Raccourci api.ini.lnk, profil adm.kamga, SRV-CPT01}
\textbf{Source~:} image NC-03, volume système,
\texttt{Users\textbackslash adm.kamga\textbackslash AppData\textbackslash
Roaming\textbackslash Microsoft\textbackslash Windows\textbackslash
Recent}.

\textbf{Contenu~:} chemin de la cible, taille et dates de la cible au moment
de l'ouverture, nom et numéro de série du volume, dates propres du
raccourci (création le 12/03/2025 à 22~h~44).

\textbf{Ce qu'il prouve~:} qu'une session du profil \texttt{adm.kamga} a
ouvert \texttt{api.ini} à cette heure, par l'interface de Windows.

\textbf{Ce qu'il ne prouve pas~:} qui était assis devant l'écran~; il ne dit
pas non plus que la clé a été lue en entier ni recopiée.
\end{boxobjet}

\begin{boxlecture}{Les trois questions au disque de SRV-CPT01}
\begin{tableaucours}{L{6.6cm}cL{3.4cm}}{%
    \tableauHeader{Proposition} &
    \tableauHeader{Cote} &
    \tableauHeader{Sources}
}
Le compte \texttt{adm.kamga} s'est connecté le 12~mars à 22~h~39 & \OUI{} &
4624 type 10 \\
La connexion venait de SRV-AD01 & \OUI{} & 4624 (adresse source) \\
Le fichier \texttt{api.ini} a été ouvert dans cette session & \OUI{} &
\texttt{.lnk}, \texttt{RecentDocs} \\
La clé a été copiée hors du serveur & \NEK{} & Archive créée, aucune trace de
transfert sur ce disque \\
Les journaux de SRV-CPT01 ont été effacés & \NON{} & Pas de 1102, journal
continu \\
M.~KAMGA lui-même était aux commandes & \NEK{} & Aucune pièce sur la
personne \\
\end{tableaucours}

\medskip
\textbf{Lecture.} Trois \OUI{}, un \NON{}, deux \NEK{}. Le disque répond
aux questions de connexion et d'ouverture~; il ne répond pas à celle du
transfert, qui relève de l'analyse du réseau (Chapitre~\ref{chap:reseau}).
\end{boxlecture}

% ============================================================================
\section{Les Heures et leurs Pièges}
\label{sec:c04-heures}
% ============================================================================

NTFS enregistre ses dates en temps universel (UTC). L'outil d'analyse les
affiche souvent dans le fuseau de la station d'examen, et les journaux
Windows dans celui que l'on choisit. Le Cameroun est à UTC+1, sans heure
d'été. Une heure lue sans son fuseau peut donc être fausse d'une heure, et
une chronologie qui mélange les deux devient incohérente.

\begin{tableaucours}{L{4.0cm}L{8.2cm}}{%
    \tableauHeader{Piège} &
    \tableauHeader{Parade}
}
Fuseau implicite & Noter le fuseau de chaque source, convertir tout en UTC
dans les notes de travail \\
Horloge déréglée & Mesurer la dérive (heure système contre heure de
référence) au moment de l'acquisition \\
Date modifiée par l'attaquant & Comparer \texttt{\$STANDARD\_INFORMATION} et
\texttt{\$FILE\_NAME}, recouper avec \texttt{\$UsnJrnl} et les journaux \\
Résolution différente & Les journaux à la seconde, NTFS à la centaine de
nanosecondes~: on compare à la précision la plus faible \\
\end{tableaucours}
\vspace{-0.8em}
\begin{center}{\small\itshape Tableau~4.3. Les pièges de l'heure}\end{center}

À NYANGA CRÉDIT, les serveurs se synchronisaient sur un serveur de temps
interne, lui-même calé sur une source publique~; la dérive mesurée le
15~mars sur SRV-CPT01 était de 2~secondes. Les heures du chapitre sont
données en heure locale (UTC+1), comme partout dans le livre.

% ============================================================================
\section{Ce que l'Analyse du Disque demande à chaque Métier}
\label{sec:c04-metiers}
% ============================================================================

\begin{boxmetier}{Analyste en réponse à incident}
\textbf{Examiner un disque}
\begin{enumerate}
    \item Écrire les questions avant d'ouvrir l'image.
    \item Pour chaque question, choisir les artefacts et vérifier qu'ils
          existent dans cette configuration.
    \item Croiser au moins deux artefacts indépendants pour chaque
          conclusion.
    \item Noter le fuseau de chaque heure.
    \item Dire, pour chaque question, ce qui a été lu et ce qui ne l'a pas
          été.
\end{enumerate}
\end{boxmetier}

\begin{boxmetier}{Application OPJ/Expert Judiciaire}
L'expert sépare ce que le disque établit sur un compte de ce qu'il ne dit
pas sur une personne. « Le compte \texttt{adm.kamga} a ouvert le fichier » est
une constatation~; « M.~KAMGA a ouvert le fichier » n'en est pas une.
\end{boxmetier}

\begin{boxmetier}{RSSI}
Activer, avant tout incident, l'audit de la création de processus et
dimensionner les journaux pour qu'ils couvrent au moins trois mois.
Sans l'événement 4688, l'enquête de NYANGA a dû se contenter d'Amcache pour
dater les exécutions.
\end{boxmetier}

\separateur

\begin{boxresume}
\textbf{Ce qu'il faut retenir de ce chapitre~:}
\begin{itemize}
    \item Un disque se lit par couches~: partitions, système de fichiers,
          fichiers, artefacts.
    \item NTFS garde deux jeux de dates (\texttt{\$STANDARD\_INFORMATION} et
          \texttt{\$FILE\_NAME})~; la date d'accès est peu fiable.
    \item Un fichier supprimé se récupère par ses métadonnées tant que
          l'enregistrement n'est pas réutilisé, sinon par ses signatures.
    \item Les artefacts Windows répondent à des questions précises~;
          ils dépendent de la configuration et se croisent.
    \item À NYANGA, le disque établit la connexion de \texttt{adm.kamga}
          depuis SRV-AD01 et l'ouverture du fichier de clé~; le transfert
          reste \NEK{} sur ce seul disque.
\end{itemize}
\end{boxresume}

\begin{boxinfo}
\textbf{Points de contrôle.} Vous avez compris ce chapitre si vous savez~:
\begin{itemize}
    \item citer l'artefact qui répond à une question donnée~;
    \item expliquer pourquoi un fichier supprimé est souvent récupérable~;
    \item convertir une heure entre UTC et l'heure du Cameroun, et dire
          pourquoi c'est nécessaire.
\end{itemize}
\end{boxinfo}

% ============================================================================
\section*{Exercices}
% ============================================================================

\begin{boxexo}[title={\ding{45}\quad Exercice~4.1. Quel artefact~? \niveaubase}]
Pour chaque question, citez l'artefact Windows qui y répond le mieux~:
(a) un compte a-t-il été créé le 9~mars sur SRV-AD01~? (b) quelqu'un s'est-il
connecté par le bureau à distance~? (c) un service inconnu a-t-il été
installé~? (d) le journal de sécurité a-t-il été vidé~? (e) quel document
Word l'utilisateur a-t-il ouvert en dernier~?
\end{boxexo}

\begin{boxexo}[title={\ding{45}\quad Exercice~4.2. Récupérer cpt.7z \niveaumoyen}]
\begin{enumerate}[(a)]
    \item Expliquez pourquoi l'archive \texttt{cpt.7z} a pu être récupérée
          avec son nom et ses dates.
    \item Si l'enregistrement MFT avait été réutilisé, quelle méthode
          resterait-il, et qu'aurait-on perdu~?
    \item L'archive est chiffrée. Que peut-on en tirer sans le mot de passe,
          et que faire pour l'obtenir~?
\end{enumerate}
\end{boxexo}

\begin{boxexo}[title={\ding{45}\quad Exercice~4.3. Deux heures pour un fichier \niveaumoyen}]
L'outil d'analyse affiche, pour \texttt{api.ini.lnk}, une création à
« 21:44:07 », sans indication de fuseau. Le journal de sécurité, exporté en
heure locale, donne l'ouverture de session à 22:39:12.
\begin{enumerate}[(a)]
    \item Le raccourci a-t-il été créé avant l'ouverture de session~?
    \item Comment lever le doute~?
    \item Réécrivez les deux heures en UTC et en heure du Cameroun.
\end{enumerate}
\end{boxexo}

\begin{boxexo}[title={\ding{45}\quad Exercice~4.4. Du compte à la personne \niveauexpert}]
Le directeur général lit l'encadré « Les trois questions au disque de
SRV-CPT01 » et demande~: « Donc
M.~KAMGA est coupable~? »
\begin{enumerate}[(a)]
    \item Répondez-lui en séparant constatation et analyse.
    \item Quelles pièces, sur ce disque ou ailleurs, pourraient faire passer
          la dernière ligne de \NEK{} à \OUI{} ou à \NON{}~?
    \item Quelle pièce du Chapitre~\ref{chap:investigation} rend déjà peu
          probable une action de M.~KAMGA depuis le siège~?
\end{enumerate}
\end{boxexo}

\bigskip

\begin{boxinfo}
\textbf{Pour aller plus loin.}
\begin{itemize}
    \item B.~Carrier, \textit{File System Forensic
          Analysis}~\cite{carrier2005}, pour NTFS et ext.
    \item The Sleuth Kit et Autopsy~\cite{sleuthkit}, outils libres
          d'analyse de disque.
\end{itemize}
\end{boxinfo}

\begin{boxinfo}
\textbf{Transition.} La session de \texttt{adm.kamga} venait du contrôleur
de domaine. Comment l'attaquant y a-t-il obtenu ces identifiants~? La
réponse est dans la mémoire de SRV-AD01~: Chapitre~\ref{chap:memoire}.
\end{boxinfo}
